\section{Stable sign data and nodal measure}
\label{sec:signs}

Fix the data of Lemma~\ref{env:gaps}, and recall that
$B=\sqrt{2+|\nabla_h f|_h^2}$. We select one of the real quasimodes
from Section~\ref{sec:four-waves} with many strict opposite-sign tests
along short coordinate segments. The score bounds nodal measure
from below, and every fixed strict lower bound for it survives
sufficiently small uniform perturbations of the function.
All limits in this section hold with the data and requested packet
accuracy fixed. The positive constants in the score and measure bounds,
however, depend only on the fixed chart and metric comparability
throughout $\mathcal U$.

\begin{definition}[Sign score]\label{score:definition}
Subdivide $D$, up to its grid boundaries, into finitely many disjoint
open coordinate cubes $Q$, fine enough that numbers $B_Q>0$ satisfy
\[
 B_Q\le B(x)\le 2B_Q\qquad (x\in Q).
\]
Choose $c_0>0$ as in Lemma~\ref{score:expectation} below and put
\[
 \ell_Q=\frac{c_0}{kB_Q},\qquad
 Q^- =\{x\in Q:\operatorname{dist}_{\infty}(x,\partial Q)>\ell_Q\},
 \qquad T_k=\sum_Q\ell_Q^{-1}|Q|.
\]
For a real continuous function $\psi$ on $M$, define
\begin{equation}\label{score:formula}
 \mathcal S(\psi)=\sum_Q\ell_Q^{-1}
 \int_{Q^-}\sum_{i=1}^4
 \mathbf 1_{\{\psi(x)\psi(x+\ell_Q d_i)<0\}}\,dx,
\end{equation}
where $d_i$ are the four coordinate unit vectors and $|Q|$ denotes
coordinate volume.  The subdivision and $k$ are part of the score's
definition and remain fixed when $\psi$ varies.
\end{definition}

In particular, $0\le\mathcal S(\psi)\le4T_k$.  If $\psi_n\to\psi$
uniformly, every strictly negative endpoint product for $\psi$ is
eventually negative for $\psi_n$.  Fatou's Lemma gives
\begin{equation}\label{score:lsc}
 \mathcal S(\psi)\le\liminf_{n\to\infty}\mathcal S(\psi_n).
\end{equation}
Multiplication by a positive continuous function, or by a nonzero real
constant, preserves every endpoint-sign comparison and hence the score.

\begin{lemma}[The score bounds nodal measure]\label{score:measure}
There is $c_{\mathrm H}>0$, independent of the subdivision, $k$, and
the envelope, such that, for every $g\in\mathcal U$ and every real
$\psi\in C(M)$,
\begin{equation}\label{score:hausdorff}
 \mathcal H_g^3(\{\psi=0\})\ge c_{\mathrm H}\mathcal S(\psi).
\end{equation}
\end{lemma}

\begin{proof}
Write $Z=\{\psi=0\}$ and, for fixed $Q,i$, let $E_{Qi}$ be the
set indicated in the integrand of \eqref{score:formula}.  This is a
Borel set.  The full segment from $x$ to $x+\ell_Qd_i$ lies in $Q$
when $x\in Q^-$.  The Intermediate Value Theorem therefore gives,
in coordinate space,
\[
 E_{Qi}\subset (Z\cap Q)-[0,\ell_Q]d_i.
\]
Cover $Z\cap Q$ by countably many Euclidean balls of radii
$r_j<\delta$.  The backward sweep of each ball is a cylinder with
rounded ends whose four-dimensional volume is at most
$C(\ell_Qr_j^3+r_j^4)$.  Countable subadditivity yields
\[
 |E_{Qi}|\le C(\ell_Q+\delta)\sum_jr_j^3.
\]
If the coordinate Hausdorff three-measure of $Z\cap Q$ is finite,
take the infimum over such covers and then let $\delta\downarrow0$.
The definition of Hausdorff measure, with the harmless dimensional
factor converting diameters to ball radii, gives
\begin{equation}\label{score:sweep}
 \ell_Q^{-1}|E_{Qi}|
 \le C\mathcal H^3_{\mathrm{Eucl}}(Z\cap Q).
\end{equation}
If this Hausdorff measure is infinite, the same inequality is immediate.
This proof uses only measurable covering cylinders and does not require
any regularity of $Z$.

To compare the two Hausdorff measures, multiply each coordinate function
by a smooth cutoff equal to one near $\overline D$ and compactly
supported in $P$, and extend by zero to $M$.  The resulting smooth map
$F:M\to\mathbb R^4$ agrees with the coordinate map near $\overline D$.
Uniform metric comparability bounds $|dF|_g$ by a fixed constant.
Integrating along curves shows that $F$ is globally Lipschitz for the
Riemannian distance of every $g\in\mathcal U$.  Thus
\[
 \mathcal H^3_{\mathrm{Eucl}}(Z\cap Q)
 \le C\mathcal H_g^3(Z\cap Q).
\]
Sum \eqref{score:sweep} over the four directions and the disjoint
open cubes.  The sets $Z\cap Q$ are disjoint Borel sets, so their
Hausdorff measures sum to at most $\mathcal H_g^3(Z)$.  This proves
\eqref{score:hausdorff}, including the case of infinite measure.
\end{proof}

\begin{lemma}[Expected score]\label{score:expectation}
One can choose $c_0>0$ and $c_{\mathrm E}>0$, using only uniform
metric comparability, such that the Gaussian superposition
\eqref{wave:random} satisfies, for all sufficiently large integers $k$,
\begin{equation}\label{score:expected}
 \mathbb E\mathcal S(u)\ge c_{\mathrm E}T_k.
\end{equation}
The threshold for $k$ may depend on all fixed data, including the envelope
and the packet accuracy, whereas $c_0,c_{\mathrm E}$ do not.
\end{lemma}

\begin{proof}
On $D$ the macroscopic cutoff is one.  Normalize the random part by
$e^{kf(x)}$ and write
\[
 v_{yj}(x)=e^{-kf(x)}p_{yj,k}(x),\qquad
 X(x)=\operatorname{Re}\sum_{y,j}(G_{yj}+iG'_{yj})v_{yj}(x).
\]
The two Gaussian coefficients of each packet are independent, so
\begin{equation}\label{score:covariance}
 \operatorname{Var}X(x)=\sum_{y,j}|v_{yj}(x)|^2,\qquad
 \operatorname{Cov}(X(x),X(x'))
 =\sum_{y,j}\operatorname{Re}\bigl(v_{yj}(x)\overline{v_{yj}(x')}\bigr).
\end{equation}
The nearest grid center is at distance $O(k^{-1})$ from $x$ and lies in
$U_3$.  The two-sided quadratic phase bounds and the amplitude estimate
in Proposition~\ref{wave:packets} show that its normalized packet has
absolute value bounded below by a positive constant.  Hence
$\operatorname{Var}X(x)\ge c>0$, uniformly on $D$ for the fixed data.

Fix $x\in Q^-$ and put $x_i=x+\ell_Qd_i$.  Call a packet near if
$|y-x|\le k^{-2/5}$.  For all other packets, both $x$ and $x_i$ are
at distance at least $\tfrac12 k^{-2/5}$ from the center for large $k$.
The Gaussian decay and the $O(k^4)$ packet count bound their total
contributions to the variances and covariances by
$Ck^4e^{-ck^{1/5}}$.  These are negligible even relative to the positive
variance lower bound.

For near packets the cutoff is one at both endpoints.  Put
$\beta_{yj}=b_j^\flat(y)$.  Along the segment,
\[
 d(\Phi_{yj}-f)=i\beta_{yj}+O(k^{-2/5}).
\]
The packet amplitudes are $1+O(|x-y|+k^{-1})$, so their endpoint
ratios tend uniformly to one.  Since $k\ell_Q=c_0/B_Q$ is bounded,
integration of the preceding differential identity gives
\begin{equation}\label{score:endpoints}
 v_{yj}(x_i)=v_{yj}(x)
 \left(\exp\!\left(\frac{ic_0\beta_{yj}(d_i)}{B_Q}\right)+o(1)\right).
\end{equation}
Every error here is uniform over the indicated points and packets for
the fixed data; smooth dependence of $b_j(y)$ on the center is not needed.

Let $S_0=\sum_{\mathrm{near}}|v_{yj}(x)|^2$.  The nearest packet gives
$S_0\ge c$, while \eqref{score:endpoints} and the tail estimate show
that the two endpoint variances are $S_0(1+o(1))$.  Thus their centered
Gaussian correlation is
\begin{equation}\label{score:correlation}
 \mathcal R_i=
 \sum_{\mathrm{near}}\frac{|v_{yj}(x)|^2}{S_0}
 \cos\!\left(\frac{c_0\beta_{yj}(d_i)}{B_Q}\right)+o(1).
\end{equation}

Uniform metric comparability, $|b_j(y)|_h=B(y)$, and continuity of the
fixed $B$ give constants $a_0,a_1>0$, independent of the envelope, with
\[
 a_0B_Q\le|\beta_{yj}|_{\mathrm{Eucl}}\le a_1B_Q
 \qquad\text{for every near packet and all sufficiently large }k.
\]
Indeed $B(x)\in[B_Q,2B_Q]$, and $B(y)/B(x)\to1$ uniformly when
$|y-x|\le k^{-2/5}$.  Choose $c_0a_1\le1$.  The elementary inequality
$1-\cos s\ge s^2/4$ for $|s|\le1$ implies
\[
 \sum_{i=1}^4\left(1-\cos\frac{c_0\beta_{yj}(d_i)}{B_Q}\right)
 \ge\frac{c_0^2|\beta_{yj}|_{\mathrm{Eucl}}^2}{4B_Q^2}
 \ge\kappa>0,
\]
where $\kappa$ is independent of the envelope. The weights in
\eqref{score:correlation} are the same for all four coordinate directions.
Weighted averaging therefore gives $\sum_i(1-\mathcal R_i)\ge\kappa/2$,
so for at least one $i$, $\mathcal R_i\le1-\kappa/8$. This bound uses
the length of each carrier covector, not the Gram determinant of the
triple chosen at its center; hence $\kappa$ is uniform over the amplified
envelopes.

For centered Gaussian variables with positive variances and correlation
$r\in[-1,1]$, the probability of opposite signs is $\arccos(r)/\pi$.
To see this, express the variables as projections of a standard planar
Gaussian onto unit vectors with angle $\arccos r$ and use rotation
invariance; the degenerate endpoints follow directly.  Therefore
\[
 \sum_{i=1}^4\mathbb P\{X(x)X(x_i)<0\}\ge q_0>0
\]
with $q_0$ independent of the envelope.

The normalized deterministic mean is
$m(x)=e^{k(f_0(x)-f(x))}\cos(k\phi(x))$, with
$|m(x)|\le e^{-k\sigma}$ on $D$ by Lemma~\ref{env:gaps}.
For a centered one-dimensional Gaussian of standard deviation $s>0$,
adding a mean $m$ changes its sign with probability at most
$\mathbb P\{|X|\le|m|\}\le C|m|/s$.
The variance lower bounds at both endpoints consequently change each
opposite-sign probability by $o(1)$, uniformly in $x$.  Thus
\[
 \sum_{i=1}^4\mathbb P\{u(x)u(x_i)<0\}\ge q_0/2
\]
for large $k$.  Tonelli's Theorem and
$|Q^-|/|Q|\to1$ for the fixed finite subdivision now prove
\eqref{score:expected} with, for example, $c_{\mathrm E}=q_0/4$.
\end{proof}

\begin{proposition}[A simultaneous good realization]\label{score:good}
For every predetermined pair of finite orders $\ell,J\ge0$, all
sufficiently large integers $k$ admit a choice of the real Gaussian
coefficients for which the smooth real function $u$ satisfies
\begin{align}
 &u=u_k^0\quad\text{outside a compact subset of }U_2,
 \qquad \supp R\Subset U_2,
 \qquad R=(\Delta_h+\Lambda_k)u,\label{score:support}\\
 &\sup_P W^{-1}\sum_{|\alpha|\le\ell}|\partial^\alpha R|
 \le C_{\ell,J}k^{-J},\label{score:residual}\\
 &W^{-1}|\partial^\alpha u|\le C_\alpha k^{|\alpha|+5}
 \quad\text{on }P\text{ for every fixed multi-index }\alpha,\notag\\
 &|(u,k^{-1}du)|\ge k^{-200}W
 \quad\text{on }\overline U_2,\label{score:bounds}\\
 &\mathcal S(u)\ge c_{\mathrm s}k\int_D B(x)\,dx.
 \label{score:large}
\end{align}
Here $c_{\mathrm s}>0$ is independent of the envelope and the accuracy
orders.  The constants $C_{\ell,J},C_\alpha$ and the frequency threshold
may depend on these fixed data, but the polynomial exponents in these
bounds do not increase with the residual accuracy.
\end{proposition}

\begin{proof}
Choose the packet accuracy in Lemma~\ref{wave:residual} to give
\eqref{score:residual}.  Let $G_k$ be the event in
Lemma~\ref{wave:noncancellation}: all coefficients have absolute value
at most $k$, and the stated noncancellation bound holds.  Its probability
tends to one, and its coefficient bound gives the support, residual,
and derivative conclusions above.  By Lemma~\ref{score:expectation}
and the deterministic upper bound $\mathcal S\le4T_k$,
\[
 \mathbb E\bigl[\mathcal S(u)\mathbf1_{G_k}\bigr]
 \ge\bigl(c_{\mathrm E}-4\mathbb P(G_k^c)\bigr)T_k
 \ge\frac{c_{\mathrm E}}2T_k
\]
for all large $k$.  Hence there is a realization in $G_k$ with
$\mathcal S(u)\ge c_{\mathrm E}T_k/4$.  This conclusion uses no
independence between the score and $G_k$.  Finally,
\[
 T_k=\frac{k}{c_0}\sum_QB_Q|Q|
 \ge\frac{k}{2c_0}\int_D B(x)\,dx,
\]
so \eqref{score:large} holds with
$c_{\mathrm s}=c_{\mathrm E}/(8c_0)$.
\end{proof}
